\documentclass[11pt]{article}

\usepackage[a4paper,margin=1in]{geometry}
\usepackage{kotex}
\usepackage{amsmath,amssymb}
\usepackage{booktabs}
\usepackage{tabularx}
\usepackage{array}
\usepackage{graphicx}
\usepackage{hyperref}
\usepackage{xcolor}

\hypersetup{
  colorlinks=true,
  linkcolor=blue,
  urlcolor=blue,
  citecolor=blue,
}

\newcommand{\EmoNet}{EmoNet}

\title{\EmoNet: 감정 동역학 기반 한국어 응답 생성 프레임워크\\\large Working Paper Draft}
\author{Author TBD}
\date{2026-04-10}

\begin{document}

\maketitle

\begin{abstract}
\EmoNet{}은 한국어 감정 응답 생성을 입력 자극 인코딩, 감정 동역학, branch/trajectory 요약, 잠재표현, 스타일 회귀, LLM 표면화로 분해하는 프레임워크다. 이 시스템의 목적은 단순히 응답 품질을 높이는 것보다, 감정 제어 실패가 어느 단계에서 발생하는지 진단 가능한 중간표현을 만드는 데 있다. 최신 cycle에서는 calibrated reference configuration을 도입해 dominant branch 평균 길이를 1.0539에서 70.4684로 늘리고 길이 1 비율을 0.9734에서 0.0948로 낮췄다. 또한 raw trajectory를 직접 읽는 heuristic readout 대신 GPT-5.4 기반 episode interpretation 경로를 추가해 trajectory 수준 affect episode를 더 풍부하게 기술할 수 있게 했다. 그러나 end-to-end generation에서는 여전히 \texttt{stim\_only}가 mean total 3.5404로 가장 높고, \texttt{episode\_trace}는 3.4673, \texttt{emonet\_full}은 3.2478에 머문다. predictor 비교에서도 learned $z \rightarrow s$ 경로는 text baseline보다 강하지 않다. 따라서 현재 \EmoNet{}의 가장 강한 기여는 최종 성능 우위 자체보다, branch collapse 완화와 trajectory interpretation을 통해 감정 처리의 내부 병목을 측정 가능한 형태로 드러냈다는 데 있다.
\end{abstract}

\section{서론}

대규모 언어모델 기반 응답 생성은 문장 자연성과 의미 보존 측면에서 빠르게 향상되었지만, 감정이 내부에서 어떻게 형성되고 어떤 경로를 거쳐 최종 말투로 번역되는지는 여전히 설명하기 어렵다. 많은 접근은 감정 라벨이나 속성 토큰을 직접 prompt에 주입하는 방식에 머무르기 때문에, 실패가 입력 인코더, 내부 동역학, 잠재표현, 스타일 회귀, 또는 최종 surface realization 중 어디에서 발생했는지 분리해 보기 어렵다.

\EmoNet{}은 이 문제를 단계별 파이프라인으로 재구성한다. 입력 텍스트는 4차원 affective stimulus로 변환되고, clustered neuro-affective dynamics core는 이 자극을 시간축으로 전개하면서 branch와 trajectory를 생성한다. 이후 지배적인 branch 요약은 잠재표현 $z$로 압축되고, 다시 스타일 벡터 $s$로 회귀되며, 최종 LLM은 이 상태를 자연어 표면으로 실현한다.

이 논문의 핵심 질문은 두 가지다. 첫째, branch-based intermediate representation이 실제로 trivial collapse에서 벗어났는가. 둘째, richer internal trajectory가 실제 generation quality 개선으로 이어지는가. 현재 결과는 첫 질문에는 비교적 긍정적으로 답하지만, 둘째 질문에는 아직 제한적인 답만을 제공한다. 이 점이 오히려 중요하다. \EmoNet{}의 현재 강점은 모든 평가에서 최고 점수를 내는 것이 아니라, 내부 구조 개선과 최종 출력 사이의 간극을 정량적으로 드러냈다는 데 있다.

\section{시스템 개요}

\EmoNet{}의 active pipeline은 다음과 같이 요약된다.

\begin{equation}
x \rightarrow E_{\mathrm{aff}}(x) \rightarrow v_{\mathrm{stim}}
\rightarrow \mathcal{N}_{\mathrm{cad}} \rightarrow H_{\mathrm{traj}}
\rightarrow z \rightarrow s \rightarrow G_{\mathrm{prompt}} \rightarrow y
\end{equation}

여기서 $x$는 입력 텍스트, $v_{\mathrm{stim}}$은 4차원 affective stimulus, $\mathcal{N}_{\mathrm{cad}}$는 clustered neuro-affective dynamics core, $H_{\mathrm{traj}}$는 tick 단위 trajectory memory, $z$는 dominant branch 기반 잠재표현, $s$는 스타일 회귀 결과, $y$는 최종 응답이다.

\paragraph{Affective stimulus encoding.}
입력 텍스트는 dopamine, serotonin, norepinephrine, melatonin의 4축 stimulus로 투영된다. 이 단계는 응답의 감정 톤을 직접 출력하지 않고, 후속 동역학을 구동하는 저차원 제어 신호를 만드는 역할을 한다.

\paragraph{Clustered affective dynamics.}
중심 코어는 256개 노드로 구성된 clustered dynamics network이며, stimulus에 따라 활성화, 억제, persistence를 시간축으로 전개한다. 각 rollout은 tick별 활성 노드와 firing edge를 trajectory로 기록한다.

\paragraph{Branch extraction and latent compression.}
trajectory에서 dominant branch를 추출하고, 이를 잠재표현 $z$로 압축한다. 본래 이 단계의 가장 큰 문제는 branch가 거의 항상 길이 1에서 끝나는 collapse 현상이었고, 이번 cycle의 상당수 실험은 이 문제를 줄이는 데 집중되었다.

\paragraph{Style regression and response realization.}
잠재표현 $z$는 40차원 스타일 표현 $s$로 회귀된다. 최종 generation 단계에서 LLM은 감정 상태를 계산하는 주체라기보다, \EmoNet{}이 제공한 stimulus, trajectory summary, style signal을 언어 표면으로 번역하는 주체로 동작한다.

\paragraph{Trajectory interpretation.}
최근 active line의 중요한 변경은 heuristic top-emotion readout 대신 raw trajectory를 GPT-5.4가 episode 단위로 해석하도록 경로를 바꿨다는 점이다. 이 해석은 stimulus label 하나를 복원하려는 시도보다, stimulus, appraisal, trajectory pattern, action tendency를 함께 읽어 ``emotion episode''를 설명하는 쪽에 가깝다.

\section{실험 설정}

\subsection{현재 기준 산출물}

현재 논문 초안은 2026-04-09 refresh bundle과 2026-04-10 episode comparison summary를 기준으로 작성한다. 핵심 숫자는 다음과 같다.

\begin{itemize}
  \item Full export rows: 51,628
  \item Branch mean length: 70.4684
  \item Branch len1 ratio: 0.0948
  \item Parsed style rows: 3,971
  \item Keep-valid rows: 1,717
  \item Rebalanced supervised rows: 1,800
  \item Encoder head validation MAE: 0.10273
  \item Decoder validation MAE: 0.117898
\end{itemize}

\subsection{평가 관점}

이번 원고의 핵심 평가는 세 층으로 나뉜다.

\begin{enumerate}
  \item \textbf{구조적 회복}: branch collapse가 완화되었는지, calibration이 reference configuration을 정당화하는지 본다.
  \item \textbf{표현력}: raw trajectory가 heuristic emotion label보다 더 풍부한 affect episode를 담는지 본다.
  \item \textbf{종단 품질}: 내부 trajectory 개선이 실제 generation score 향상으로 번역되는지 본다.
\end{enumerate}

\subsection{비교 조건}

현재 generation 비교 조건은 다음 네 묶음으로 정리할 수 있다.

\begin{itemize}
  \item Baseline: \texttt{stim\_only}, \texttt{direct}
  \item Trace-conditioned: \texttt{raw\_trace}, \texttt{hybrid\_trace}, \texttt{appraisal\_trace}
  \item Episode-conditioned: \texttt{episode\_trace}, \texttt{hybrid\_episode}
  \item Full path: \texttt{emonet\_full}
\end{itemize}

이들 비교는 ``내부 상태가 더 풍부해질수록 최종 응답도 좋아지는가''를 직접 확인하기 위한 것이다.

\subsection{이 원고의 평가 원칙}

현재 단계에서 중요한 원칙은 과장하지 않는 것이다. 본문은 다음 세 가지를 함께 보여준다.

\begin{itemize}
  \item branch collapse recovery는 강하게 주장할 수 있다.
  \item trajectory interpretation richness는 조심스럽게 주장할 수 있다.
  \item end-to-end superiority는 아직 주장하면 안 된다.
\end{itemize}

\section{주요 결과}

\subsection{Branch collapse는 크게 완화되었다}

가장 먼저 확인할 수 있는 결과는 dominant branch의 trivial collapse가 구조적으로 줄었다는 점이다. Table~\ref{tab:branch-recovery}는 structural recovery 전후의 branch length 분포를 요약한다.

\begin{table}[t]
\centering
\caption{Structural recovery 전후 dominant branch length 분포.}
\label{tab:branch-recovery}
\begin{tabular}{lrrrrrr}
\toprule
Setting & Rows & Mean & Len1 ratio & p50 & p90 & Max \\
\midrule
Before structfix & 51,628 & 1.0539 & 0.9734 & 1 & 1 & 8 \\
After structfix & 51,628 & 70.4684 & 0.0948 & 69 & 126 & 126 \\
\bottomrule
\end{tabular}
\end{table}


평균 branch 길이는 1.0539에서 70.4684로 늘었고, 길이 1 branch 비율은 0.9734에서 0.0948로 감소했다. 이는 branch-based intermediate representation이 더 이상 거의 항상 1-step에서 끝나지 않는다는 뜻이다. 다만 상위 tail은 여전히 ceiling에 붙어 있으며, after setting에서도 $p90 = p95 = \max = 126$이다. 즉 collapse는 크게 줄었지만 upper-tail saturation은 아직 남아 있다.

\subsection{Reference configuration은 calibration으로 뒷받침된다}

현재 reference configuration은 임의 초기값이 아니라 calibration experiment를 통해 선택되었다. 결합 검증 기준에서 \texttt{no\_activity\_ratio = 0.05}, \texttt{len1\_ratio = 0.05}, \texttt{hit\_max\_ticks\_ratio = 0.25}, \texttt{mean\_first\_active\_tick = 2.40}, \texttt{mean\_branch\_len = 79.78}이 확인되었다. 이 결과는 현재 설정이 dead sample 감소, branch collapse 완화, 과도한 persistence 억제를 동시에 고려한 선택이라는 점을 보여준다.

\subsection{Style target bias는 여전히 강하다}

style supervision은 consistency 측면에서는 정리되어 있지만, 분포 자체는 지나치게 safe하고 cooperative하다. keep-valid subset 평균에서 softness는 0.9276, calmness는 0.9132, cooperativeness는 0.9202, positivity는 0.9051이다. 반면 hostility와 resentment는 각각 0.0003, despair는 0.0044, volatility는 0.0022에 불과하다. 이는 내부 affect signal이 어느 정도 회복되어도, supervision target 자체가 최종 surface를 다시 순화시키는 강한 편향으로 작동할 수 있음을 시사한다.

\subsection{Predictor는 아직 text baseline을 넘지 못한다}

Table~\ref{tab:predictor-comparison}은 predictor 성능을 비교한다.

\begin{table}[t]
\centering
\caption{Predictor mean MAE 비교. 낮을수록 좋다.}
\label{tab:predictor-comparison}
\begin{tabular}{lrr}
\toprule
Predictor & Mean MAE & Gain vs baseline \\
\midrule
Mean baseline & 0.117288 & 0.000000 \\
Stim-only ridge & 0.116513 & 0.000775 \\
Text tf-idf ridge & 0.114628 & 0.002660 \\
Legacy z64 ridge & 0.116846 & 0.000442 \\
Structfix learned z64 ridge & 0.117328 & -0.000040 \\
\bottomrule
\end{tabular}
\end{table}


text tf-idf ridge는 mean MAE 0.114628로 가장 낮고, structfix learned z64 ridge는 0.117328로 baseline보다도 소폭 나쁘다. 이는 branch dynamics recovery가 곧바로 better $z \rightarrow s$ prediction으로 이어지지 않았음을 의미한다. 현재 병목은 dynamics 자체라기보다, trajectory 정보를 latent와 style path로 압축하는 과정에 남아 있을 가능성이 높다.

\subsection{Generation에서는 richer trace가 naive \texttt{emonet\_full}보다 낫지만 최고는 아니다}

최신 GPT-5.4 judge summary 기준 generation 결과는 Table~\ref{tab:generation-comparison}과 같다.

\begin{table}[t]
\centering
\caption{Latest episode-v2 generation summary scored by GPT-5.4 judge.}
\label{tab:generation-comparison}
\resizebox{\linewidth}{!}{%
\begin{tabular}{lrrrrrr}
\toprule
Condition & Content fit & Emotional approp. & Style match & Naturalness & Overall quality & Mean total \\
\midrule
stim\_only & 3.9123 & 3.6228 & 2.2544 & 4.4561 & 3.4561 & 3.5404 \\
direct & 3.9292 & 3.5841 & 2.1947 & 4.4159 & 3.4336 & 3.5115 \\
episode\_trace & 4.0531 & 3.5664 & 2.1327 & 4.1947 & 3.3894 & 3.4673 \\
raw\_trace & 3.9558 & 3.5133 & 1.9381 & 4.1770 & 3.2743 & 3.3717 \\
hybrid\_trace & 3.6964 & 3.5000 & 2.1964 & 3.9732 & 3.2679 & 3.3268 \\
emonet\_full & 3.6814 & 3.3805 & 2.2920 & 3.7345 & 3.1504 & 3.2478 \\
appraisal\_trace & 3.8091 & 3.3364 & 1.9727 & 3.8727 & 3.1818 & 3.2345 \\
hybrid\_episode & 3.6937 & 3.4054 & 2.3333 & 3.5045 & 3.0721 & 3.2018 \\
\bottomrule
\end{tabular}%
}
\end{table}


\texttt{episode\_trace}는 mean total 3.4673으로 \texttt{emonet\_full}의 3.2478보다 높다. \texttt{raw\_trace} 역시 3.3717로 naive \texttt{emonet\_full}보다 낫다. 이는 raw trajectory나 episode-conditioned path가 naive end-to-end path보다 유의미한 정보를 전달하고 있음을 시사한다. 그러나 최고 점수는 여전히 \texttt{stim\_only}의 3.5404이며, \texttt{direct}도 3.5115로 더 높다.

따라서 현재 결과를 가장 정직하게 해석하면 다음과 같다. internal trace quality improvement는 분명 존재하지만, 그것이 아직 final generation superiority로 안정적으로 번역되지는 않는다. 특히 \texttt{emonet\_full}은 style match가 2.2920으로 비교적 높지만 naturalness는 3.7345로 낮다. 즉 richer internal signal이 오히려 무거운 surface conditioning을 통해 자연스러움을 깎는 부작용을 내고 있을 가능성이 있다.

\subsection{Trajectory interpretation은 reportability를 높인다}

현재 raw trajectory에서 heuristic top-emotion readout만 쓰면 긍정적 고각성이나 suspended anticipation을 적절히 설명하지 못하는 사례가 남아 있다. GPT-5.4 기반 episode interpretation은 이 trajectory를 ``하나의 라벨''이 아니라 자극, appraisal, 지속성, action tendency를 포함한 episode로 요약한다. 현재 단계에서 이것을 ground-truth emotion recovery라고 주장할 수는 없지만, trajectory의 semantic reportability를 높였다는 해석은 충분히 방어 가능하다.

\section{한계와 다음 단계}

현재 \EmoNet{} 원고는 분명한 한계를 함께 가져가야 한다.

\paragraph{첫째, 최고 성능 모델이 아니다.}
현재 generation 결과에서 최고 점수는 \texttt{stim\_only}다. 따라서 이 논문은 ``best model paper''가 아니라 ``diagnostic systems paper''로 쓰는 편이 맞다.

\paragraph{둘째, style target bias가 여전히 매우 강하다.}
현재 supervision target은 calm, soft, cooperative surface에 치우쳐 있어 raw affect를 충분히 유지하기 어렵다. felt state와 response style을 분리하는 데이터 설계가 필요하다.

\paragraph{셋째, learned latent는 predictor superiority를 확보하지 못했다.}
$z \rightarrow s$ 경로는 아직 text baseline보다 강하지 않다. 이후 단계에서는 branch extractor나 dynamics보다도 latent compression과 decoder 구조를 먼저 다시 볼 필요가 있다.

\paragraph{넷째, saturation tail이 남아 있다.}
branch collapse는 완화되었지만 ceiling-near persistence는 아직 완전히 사라지지 않았다. 이는 later-phase differentiation을 평평하게 만들 수 있다.

\paragraph{다섯째, paper artifact source를 단일화할 필요가 있다.}
현재 refresh bundle과 episode-v2 summary가 서로 다른 generation slice를 보고 있기 때문에, 최종 투고 전에는 단일 run 기준으로 표와 그림을 다시 맞춰야 한다.

\section{결론}

\EmoNet{}의 current cycle이 보여주는 가장 강한 사실은 세 가지다. 첫째, branch collapse는 calibration과 structural recovery를 통해 실질적으로 완화되었다. 둘째, raw trajectory는 heuristic emotion label보다 richer affect episode를 담고 있으며, episode interpretation은 그 정보를 더 사람이 읽을 수 있는 수준으로 드러낸다. 셋째, 그럼에도 불구하고 현재 내부 품질 향상은 아직 최종 generation superiority로 완전히 연결되지 않는다.

따라서 이 원고의 가장 적절한 결론은 ``\EmoNet{}이 모든 평가에서 최강이다''가 아니라 ``감정 처리의 내부 병목을 이전보다 훨씬 더 명확하게 측정하고 분해할 수 있게 되었다''이다. 다음 단계의 우선순위는 style target bias 완화, $z \rightarrow s$ predictor 재설계, 그리고 episode-to-response surface translation 경량화다.


\section*{References}
현재 원고는 구조와 본문 정리에 집중한 working draft이며, 관련 연구 인용과 참고문헌은 다음 정리 단계에서 보강한다.

\end{document}
